\documentclass{article}
\usepackage{graphicx}
\usepackage[a4paper, left=2cm, right=2cm, top=3cm, bottom=3cm]{geometry}

\usepackage[utf8]{inputenc}
\usepackage[greek,english]{babel}
\usepackage{alphabeta}
\usepackage{amsmath}
\usepackage{caption}
\usepackage{subcaption}
\usepackage{xcolor}
\usepackage{mathabx}
\usepackage{amssymb}
\usepackage{enumitem}

\captionsetup{labelformat=empty}
\subcaptionsetup{labelformat=empty}

\title{\textbf{Προσαρμοστικός Έλεγχος - Εργαστήριο: Coupled DC Motors}}
\author{\textbf{Μέλη Ομάδας}:\\Βραχωρίτη Αλεξάνδρα (A.M.: 1092793)\\ Μπαθρέλλος Μιχαήλ (Α.Μ.: 1092665)\\Σταναθιώτη Παναγιώτα (Α.Μ.: 1092673)}
\date{20 Νοεμβρίου 2025}

\begin{document}

\maketitle 
    \section{Εργαστήριο I.1}
    \subsection{Αναγνώριση και ανάλυση του συστήματος}
    Στόχος αυτού του εργαστηρίου είναι η προσέγγιση του συστήματος με μια (κανονικοποιημένη) συνάρτηση μεταφοράς πρώτης τάξης:
    
    \begin{equation*}
        G(s) = \frac{K}{τs+1}
    \end{equation*}
    
    διαλέγοντας ως ονομαστική λειτουργία: 
    \begin{itemize}
        \item \(PWM_{nom} = 200\)
        \item \(rpm_{nom} = 19.5\)
        \item \(error_{nom} = 2\)
    \end{itemize}
    
    \hfill\break
    Θα χρειαστεί, λοιπόν, να υπολογίσουμε το κέρδος και την χρονική σταθερά του συστήματος, επομένως πρέπει να δοκιμάσουμε διάφορες τιμές PWM (θετικές και αρνητικές) στο σύστημα και να καταγράψουμε τις στροφές (rpm) του κινητήρα, αλλά και τον χρόνο της απόκρισης από το 0\% μέχρι το 63\% της μέγιστης τιμής της εξόδου. Παρακάτω παρουσιάζονται οι πίνακες των μετρήσεων.

    \begin{table}[!h]
        \centering
        \begin{tabular}{|c|c|c|c|c|}
        \hline
            \(V_{in}\) (Volt) & \(PWM_{in}\) & \(\Omega_{out}\) (rpm) & \(\tau\) (sec) & \(Κ = \frac{ΔΩ/ΔΩ_{max}}{Δu/Δu_{max}}\)\\ \hline
            2.75 & 140 &  5.4 & 0.3 & 6.46\\ \hline
            2.94 & 150 &  8.5 & 0.4 & 6.05\\ \hline
            3.14 & 160 & 11.1 & 0.4 & 5.78\\ \hline
            3.33 & 170 & 11.7 & 0.4 & 7.15\\ \hline
            3.53 & 180 & 15.4 & 0.3 & 5.64\\ \hline
            3.73 & 190 & 16.8 & 0.4 & 7.43\\ \hline
            3.92 & 200 & 19.5 &   0 & -   \\ \hline
            4.12 & 210 &   20 & 0.4 & 1.38\\ \hline
            4.31 & 220 & 21.4 & 0.5 & 2.61\\ \hline
            4.51 & 230 & 22.7 & 0.5 & 2.93\\ \hline
            4.71 & 240 &   24 & 0.4 & 3.09\\ \hline
            5.00 & 255 & 27.4 & 0.3 & 3.95\\ \hline
        \end{tabular}
        \caption{\textbf{Πίνακας 1.} Για τις παραπάνω μετρήσεις η εξωτερική τροφοδοσία του συστήματος κυμαινόταν από \(2.75\,V\) έως \(5\,V\), η αντίσταση του ποτενσιόμετρου ήταν σταθερή \(R = 5\,k\Omega\) και επιλέξαμε nominal PWM \(= 200\), δίνοντας στο σύστημά μας μόνο θετικές τιμές PWM.}
    \end{table}

    \begin{table}[!h]
        \centering
        \begin{tabular}{|c|c|c|c|c|}
        \hline
            \(V_{in}\) (Volt) & \(PWM_{in}\) & \(\Omega_{out}\) (rpm) & \(\tau\) (sec) & \(Κ = \frac{ΔΩ/ΔΩ_{max}}{Δu/Δu_{max}}\)\\ \hline
            -2.75 & -140 &  6.5 & 0.5 & 5.27\\ \hline
            -2.94 & -150 &  9.4 & 0.3 & 4.73\\ \hline
            -3.14 & -160 & 11.4 & 0.3 & 4.54\\ \hline
            -3.33 & -170 & 13.4 & 0.3 & 4.22\\ \hline
            -3.53 & -180 & 14.8 & 0.3 & 4.40\\ \hline
            -3.73 & -190 & 16.3 & 0.4 & 4.68\\ \hline
            -3.92 & -200 &   18 &   0 & -   \\ \hline
            -4.12 & -210 & 19.7 & 0.5 & 4.68\\ \hline
            -4.31 & -220 & 20.6 & 0.5 & 3.58\\ \hline
            -4.51 & -230 & 21.4 & 0.4 & 3.12\\ \hline
            -4.71 & -240 & 22.6 & 0.4 & 3.16\\ \hline
            -5.00 & -255 &   26 & 0.3 & 4.00\\ \hline
        \end{tabular}
        \caption{\textbf{Πίνακας 2.} Για τις παραπάνω μετρήσεις η εξωτερική τροφοδοσία του συστήματος κυμαινόταν από \(2.75\,V\) έως \(5\,V\), η αντίσταση του ποτενσιόμετρου ήταν σταθερή \(R = 5\,k\Omega\) και επιλέξαμε nominal PWM \(= -200\), δίνοντας στο σύστημά μας μόνο αρνητικές τιμές PWM.}
    \end{table}

    \hfill\break
    Τελικά υπολογίστηκε:
    \begin{equation*}
        K = 4.49
    \end{equation*}

    \begin{equation*}
        \tau = 0.39 \Rightarrow pole = -2.59
    \end{equation*}

    \hfill\break
    άρα η κανονικοποιημένη συνάρτηση μεταφοράς που προσεγγίζει το σύστημα είναι:
    
    \begin{equation*}
        \boxed{G(s) = \frac{11.513}{s + 2.59}}
    \end{equation*}
    
    \hfill\break
    Παρακάτω θα δούμε την απόκριση του συστήματος για κάποιες εισόδους, ώστε να επιβεβαιώσουμε ότι το αποτέλεσμά μας είναι ορθό.

    % EIKONES

    \hfill\break
    Βάσει των παραπάνω, παρατηρούμε πως το σύστημα λειτουργεί σαν βαθυπερατό φίλτρο, επιβεβαιώνοντας την ανάλυσή μας.
    
    \subsection{Υλοποίηση P και PI ελεγκτών}
    Παρακάτω παρουσιάζεται τα διάγραμμα του ελέγχου του συστήματος με ελεγκτή \(P\).

    % εικονα πχ απο 20 σε 15 rpm (η λεζαντα να λεει Kp=...)

    \hfill\break
    Παρατηρούμε πως υπάρχει μόνιμο σφάλμα. Αυτό διορθώνεται με την πρόσθεση του \(I\) όρου.

    % εικονα για PI πχ απο 20 σε 15 rpm. (Kp=... , Ki=...)

    \hfill\break
    Παρατηρούμε πως ο PI ελεγκτής απαλείφει το σφάλμα μόνιμης κατάστασης, άρα μπορούμε να ελέγξουμε ικανοποιητικά το σύστημα. Παρακάτω θα δούμε το πρόβλημα του PI ελεγκτή.
    
    \subsection{Επίδραση του ποτενσιόμετρου στη δυναμική του συστήματος}
    Η συνάρτηση μεταφοράς του συτήματος προφανώς υπολογίστηκε για σταθερή αντίσταση του ποτενσιόμετρου \(R=5kΩ\). Επομένως το tuning του PI είναι προσαρμοσμένο σε αυτό το συγκεκριμένο σύστημα. Άρα μεταβάλοντας την αντίσταση του, το σύστημα θα αλλάξει και ο PI ελεγκτής δεν θα δουλεύει πια. Αυτό φαίνεται στα παρακάτω διαγράμματα:

    % διαγραμμα με PI ελεγτκη που δεν δουλευει λογω αλλαγης αντιστασης

    Η αντίσταση του ποτενσιόμετρου προσομοιάζει παραμετρική αβεβαιότητα. Στα επόμενα εργαστήρια θα μελετήσουμε προσαρμοστικούς ελεγκτές οι οποίοι αντιμετωπίζουν αυτού του είδους αβεβαιότητες.
    
    \section{Εργαστήριο I.2}
    \subsection{Έλεγχος ταχύτητας με τη μέθοδο του κανόνα MIT}

    \section{Εργαστήριο I.3}
    \subsection{Έλεγχος ταχύτητας με τη μέθοδο MRAC (βαθμωτή περίπτωση)}

    \section{Εργαστήριο I.4}
    \subsection{Έλεγχος ταχύτητας με τη μέθοδο ADI (βαθμωτή περίπτωση)}
    
    \section{Εργαστήριο I.5}
    \subsection{}
    
    \section{Εργαστήριο I.6}
    \subsection{}
    
\end{document}
 